
\section*{Results}

  \subsection*{A modelling pipeline to test the emergence of predictions}

    We trained single-layer gated recurrent neural networks (RNNs) to denoise noisy observations $y_t$ of sequences of tokens $x_t$. The tokens were derived from MIDI files of Bach compositions scraped from the internet. Each MIDI file was first transformed into a polyphonic piano-roll matrix. Then, the matrices were tokenised into segments that shared the same active notes; i.e., a new token was emitted only when the set of sounding pitches changed. This resulted in an event-based representation that had no notion of time. The tokens were then collapsed into a 12-dimensional chromatic space to reduce their dimensionality, so that each latent token $x_t \in \{0,1\}^{12}$. We used a total of 700 unique compositions (1432 MIDI files, as some compositions were split into parts) that were randomly split into training (490 compositions), validation (70 compositions), and testing (140 compositions).

    Observations $y_t$ were then generated by corrupting the latents with Gaussian noise $y_t = x_t + \sigma\varepsilon_t$ with $\varepsilon_t \sim \mathcal{N}\left(0, 1\right)$ (Figure~\ref{fig:design}A and~B). We used 14 noise amplitudes, ranging from essentially noiseless ($\sigma = 0.001$) to a signal-to-noise ratio of 0.5 ($\sigma = 2$), to measure a potential dependence of the emergence of prediction on the noise regime.

    RNNs were first trained on denoising (a metaphor for perception) with a binary cross-entropy (BCE) loss. RNNs received as inputs sequences of observations $y_{1:T}$ and were trained to recover, through a linear readout, the latent tokens $x_{1:T}$ (Figure~\ref{fig:design}A). After training was completed, we froze the weights of the RNNs and trained a second linear readout on prediction; i.e., the RNNs were presented with the same observations $y_{1:T-1}$, but the linear readout was trained on predicting, from the hidden states at $t$, the subsequent latent $x_{t+1}$. We trained RNNs of different sizes ($N \in \{8, 16, 32, 64, 128, 256\}$ hidden units) to measure a potential dependence of the emergence of prediction on the capacity of the network.

    \begin{figure}[tb]
      \centering
      \includegraphics[width=\textwidth]{./graph/fig01_task-and-decoding.pdf}
      \caption{\textbf{Experimental design.}
      \textbf{A)}~Training pipeline. Latent token sequences $x_{1:T}$ are corrupted with Gaussian noise to produce the observations $y_{1:T}$. RNNs are first trained on denoising the observations (training fits the black connections). Afterwards, the remaining weights are frozen and a second linear readout is trained on prediction (training fits the blue connections only).
      \textbf{B)}~The first 120-step excerpt of the stimulus sequence for BWV~111 (test set). The latent tokens $x_t$ are outlined in blue, their corresponding observations $y_t$ ($\sigma=0.25$) are shown in a greyscale heatmap (darker tones for higher values).
      \textbf{C)}~Denoised output of one of the RNNs ($N=64$) for the same excerpt. The network estimates are now shown in greyscale, blue outlines mark the ground-truth latents.
      \textbf{D)}~Decoding error for the same RNN and noise regime measured with binary cross-entropy (BCE) over the 140 test compositions. The RNN denoised the input significantly better than a memory-less logistic regression; thin lines join points from the same composition. Note that BCE is a measure of error: the lower the BCE, the better the performance.}
      \label{fig:design}
    \end{figure}

    In addition to the RNNs, we fitted five benchmark models used to contextualise the performance of the target RNNs. First, to bound the denoising problem, we fitted a memory-less logistic regression mapping a single observation $y_t$ onto the concurrent latent token $x_t$. The remaining four benchmark models bound the prediction problem: (1)~a marginal predictor that systematically predicts the global prior independently of the input, quantifying the optimal chance level under BCE; (2)~a memory-less logistic regression predicting the next token $x_{t+1}$ from a single observation $y_t$, i.e., a first-order Markovian analogue of the RNNs' readout without access to the hidden states, which we term the \emph{liberal} Markovian benchmark because it sees exactly the same input as the RNNs; (3)~an alternative version of (2) where the model was fed $x_t$ instead of the noisy observation, analogous to a linear readout of the RNNs on a perfectly denoised input, which we term the \emph{conservative} Markovian benchmark; (4)~an RNN of the same class and size as the largest networks of the sweep ($N = 256$), trained end-to-end on prediction, quantifying the best prediction performance achievable within this model class from noisy inputs.

    We expected that an RNN in which prediction emerges would show a prediction BCE below that of benchmark~(3), and that an RNN in which prediction does not emerge would show a prediction BCE between those of~(3) and~(2). In neither case did we expect the RNN to reach benchmark~(4), the empirical performance ceiling.

  \subsection*{RNNs outperform memory-less models at perception}

    RNNs successfully learned to denoise the observations (Figure~\ref{fig:design}C; Supplementary Figure~\ref{figS:decoding}). In the 64-unit RNN trained on $\sigma = 0.25$ (case example used throughout the manuscript), the network (mean $\mathrm{BCE} = 0.023$) outperformed the memory-less logistic regression (mean $\mathrm{BCE} = 0.040$; Figure~\ref{fig:design}D; Cohen's $d = 5.01$, $p < 10^{-20}$), indicating that the network exploited its memory trace during denoising.

    RNNs systematically outperformed the memory-less logistic regression benchmark at network sizes $N \geq 16$ across noise levels (Supplementary Figure~\ref{figS:decoding}; all $p < 10^{-22}$, with effect sizes ranging from $d = 1.03$ to $d = 17.53$). The 8-unit networks were, however, outperformed by the logistic regression for $0.15 \leq \sigma \leq 0.2$, as expected from a hidden state narrower than the $12$-dimensional observation it must represent. For $N \geq 16$, the advantage in raw BCE was negligible in noise regimes $\sigma \leq 0.1$, increased steadily with the noise level up to $\sigma = 0.8$, and monotonically decreased thereafter. At $\sigma = 0.8$, the denoising advantage increased monotonically with $N$ up to $N = 128$, and declined slightly at $N = 256$.

    These results are consistent with RNNs exploiting memory to improve their denoising estimates. Although, strictly speaking, RNNs can approximate non-linear mappings that could theoretically outperform the memory-less logistic regression, the growth of the advantage with $\sigma$ up to $\sigma = 0.8$ favours the memory interpretation.

  \subsection*{Predictive information emerges without a predictive objective}

    The predictive readout of the case-example network ($N = 64$, $\sigma = 0.25$) outperformed both Markovian benchmark models (Figure~\ref{fig:predictions}; Cohen's $d = 0.45$, $p < 10^{-20}$ in the comparison with the conservative Markovian benchmark, which uses the ground-truth $x_t$ as input). Therefore, at least at this noise regime, the denoising-trained RNNs encoded predictive information about the upcoming token that was not directly deducible from the previous, denoised, input. As expected, however, the network did not reach the empirical performance ceiling and was outperformed by the prediction-trained RNN ($d = -0.64$, $p < 10^{-20}$). The results were consistent across the compositions of the corpus (Figure~\ref{fig:predictions}C).

    \begin{figure}[bt]
      \centering
      \includegraphics[width=\textwidth]{./graph/fig02_prediction-one-net.pdf}
      \caption{\textbf{Predictive information emerges in an RNN trained on perception.}
      \textbf{A)}~Case-example RNN ($\sigma = 0.25$, $N = 64$) predictions (greyscale heatmap) for the same excerpt shown in Figure~\ref{fig:design}; ground-truth latents $x_{t+1}$ are outlined in blue.
      \textbf{B)}~The case-example RNN (white) shows a test-set prediction BCE consistent with the internal encoding of predictions. The network BCE is compared against four benchmark models (from left to right): (yellow) RNN fully trained on prediction, the empirical performance ceiling; (green) conservative Markovian benchmark, a first-order model predicting $x_{t+1}$ from the ground-truth $x_t$, a conservative threshold for a model without predictive information; (orange) liberal Markovian benchmark, a first-order model predicting $x_{t+1}$ from the same input as the RNN, $y_t$, a liberal threshold for a model without predictive information; and (pink) global predictive distribution, the optimal chance-level model under BCE.
      \textbf{C)}~Performance is consistent across compositions. The BCE of the case-example RNN ($x$-axis) is compared against the BCE of the best three benchmark models ($y$-axis) for each composition. Points above the diagonal are compositions in which the RNN outperforms the benchmark.}
      \label{fig:predictions}
    \end{figure}

  \subsection*{Predictive information emerges only at intermediate noise regimes}

    The emergence of predictive information depended on the noise level under which the RNNs were trained (Figure~\ref{fig:noiseregime}). The prediction performance of the 64-unit network (blue line) remained relatively stable for $\sigma \leq 0.1$ but improved abruptly at $\sigma = 0.15$, despite operating on noisier inputs $y_t$. As the noise level increased further, the performance of the RNN started to decline. The denoising-trained RNN performed comparably to the conservative Markovian benchmark (green line) for $\sigma \leq 0.1$, outperformed it for $0.15 \leq \sigma \leq 0.4$, and fell behind it thereafter. However, this decline was mirrored by the prediction-trained RNN (yellow), whose performance converged with that of the denoising-trained RNN for $\sigma \geq 0.6$. Crucially, the denoising-trained RNN performed systematically better than the liberal Markovian benchmark, which receives the exact same inputs (orange line), for $\sigma \geq 0.1$. Together, these results indicate that predictive information emerges in the denoising-trained RNNs at some point between $\sigma = 0.1$ and $\sigma = 0.15$. This result parallels the pattern observed in the denoising advantage of the RNNs over the memory-less regression, which becomes substantial between $\sigma = 0.15$ and $\sigma = 0.25$ (Supplementary Figure~\ref{figS:decoding}). Results were consistent across RNNs with $N \geq 32$ (Figure~\ref{fig:sweep}).

    \begin{figure}[htb]
      \centering
      \includegraphics[width=\textwidth]{./graph/fig03_prediction-vs-noise.pdf}
      \caption{\textbf{Predictive information emerges at moderate noise levels.}
      Prediction BCE for the 64-unit denoising-trained RNN (blue) in comparison to the benchmark models for different levels of noise ($x$-axis). At $\sigma \leq 0.1$, the RNNs show similar performance to the two Markovian models (orange and green), consistent with a lack of predictive information. At $0.15 \leq \sigma \leq 0.4$, the RNN systematically outperforms both Markovian models, demonstrating the encoding of predictive information. At $\sigma = 0.6$, the performance of the denoising-trained RNN, the prediction-trained RNN (yellow; empirical performance ceiling), and the conservative Markovian benchmark (green) merge. At $\sigma > 0.6$, both RNNs are outperformed by the conservative Markovian benchmark, which receives as inputs the noise-less latents $x_t$ and thus presents the same performance at all noise levels. Neither of the RNNs, however, is outperformed by the liberal Markovian benchmark (orange), which predicts the next token based on the noisy observations $y_t$, and asymptotically approaches the performance of the chance-level model (pink) at $\sigma=2$. This pattern of results demonstrates that the denoising-trained RNN encodes predictive information at noise regimes $0.15 \leq \sigma \leq 0.4$.}
      \label{fig:noiseregime}
    \end{figure}

    \begin{figure}[p]
      \centering
      \includegraphics[width=\textwidth]{./graph/fig04_prediction-sweep.pdf}
      \caption{\textbf{Emergent prediction across noise amplitude and network size.}
      Cohen's $d$ for the contrast between each benchmark and the denoising-trained RNN (benchmark
      minus RNN, so warm colours mark the network's advantage) across the 14 noise amplitudes
      ($x$-axis) and the 6 network sizes ($y$-axis), computed over the 140 test compositions after
      averaging each composition across the five runs.
      \textbf{A)}~Against the liberal Markovian benchmark, which receives the same noisy
      observations $y_t$ as the RNN.
      \textbf{B)}~Against the conservative Markovian benchmark, which receives the noise-free
      latents $x_t$.
      \textbf{C)}~Against the prediction-trained RNN, the empirical performance ceiling the probed
      readout is not expected to reach.
      Values enclosed in parentheses mark cells that do not survive Bonferroni correction over the
      84 cells of the design.}
      \label{fig:sweep}
    \end{figure}

  \subsection*{Activation elicited by the sensory input scales with prediction error}

    We next enquired whether the networks, if they were implemented as neural populations in the brain, would elicit responses (e.g., event-related potentials) proportional to prediction error. We measured prediction error as $e_t = x_{t} - \hat{x}_{t|t-1}$, where $\hat{x}_{t|t-1}$ is the output of the linear readout trained on prediction over the denoising-trained RNN. We assumed responses would be proportional to the magnitude of the change in activation of the RNN elicited by the new observation, $\Delta h_t = h_{t} - h_{t-1}$, where $h_t$ are the activations of the hidden units after having received $y_t$. All correlations below are computed between the squared magnitudes $|e_t|^2$ and $|\Delta h_t|^2$ across the time steps of a composition.

    The update of the case-example network ($N = 64$, $\sigma = 0.25$) closely matched prediction error in the example composition (Figure~\ref{fig:pe}A, blue), with a correlation $\rho = 0.85$, $p < 10^{-20}$. However, these correlations may be partially caused by the changes in the latent $x_t$, $|x_t - x_{t-1}|^2$, which also correlate with state updates (Figure~\ref{fig:pe}A, gray). We therefore ran a partial correlation between state update and prediction error, controlled for the change in the latent. Partial correlations were strictly above zero for all the compositions in the test set (Figure~\ref{fig:pe}B; Cohen's $d = 3.59$, $p < 10^{-20}$). This result demonstrates that the RNN dynamics would have yielded neural responses that could be interpreted as prediction error.

    \begin{figure}[p]
      \centering
      \includegraphics[width=\textwidth]{./graph/fig05_pe-state-updating.pdf}
      \caption{\textbf{State updates correlate with prediction error.}
      \textbf{A)}~Correlation between state updates and prediction error for the case-example network ($\sigma = 0.25$, $N = 64$) for a single composition (BWV~111). Although state updates are strongly correlated with prediction error (blue), they are also correlated with the update of the latent $x_t-x_{t-1}$ underlying the observations (gray; note that the latent update is discrete because the latents are binary).
      \textbf{B)}~Partial correlation between prediction error and state update controlling for the latent update across all compositions of the test set for the case-example network.
      \textbf{C)}~Mean partial correlation and effect sizes across the test set for all network sizes ($y$-axis) and noise regimes ($x$-axis). Values of Cohen's $d$ enclosed in parentheses indicate contrasts that were not significant against the null hypothesis $d = 0$ after Bonferroni correction for 84 comparisons.}
      \label{fig:pe}
    \end{figure}

    This was not the case at all noise regimes. In the 64-unit networks the partial correlation rises sharply at $\sigma = 0.25$, peaks at $\sigma = 0.4$ ($\rho = 0.29$), and decays thereafter, falling below $0.1$ by $\sigma = 1.2$ (Figure~\ref{fig:pe}C). This pattern is expected from the previous results. At low noise levels, the observation can simply be passed forward by the networks, and the state update merely reflects this pass-forward computation. On the other hand, when the observation is dominated by noise, the network update tracks neither the stimulus nor the prediction error. It is at intermediate noise levels, at which the networks leverage their predictions, that state updates reflect prediction error.

    A similar profile is apparent at other network sizes $N$, although the noise regime that promotes the encoding of prediction error widens as network size $N$ increases. This effect is probably driven by the dependence of the control term of the partial correlation on network capacity: all networks with $N \geq 16$ solve the denoising task about equally well (test BCE $0.10$--$0.12$ at $\sigma = 0.4$), but they do not use their units the same way. A network with a larger state space can afford to carry a redundant, stimulus-locked copy of the current latent token alongside whatever temporal integration the noise demands. The 8-unit networks, with a capacity lower than the dimensionality of the inputs, show a pattern of results difficult to interpret, with substantial partial correlations at all noise levels. A possible interpretation of these results is that the correlations reflect a $12 \rightarrow 8$ compression strategy of the networks that is not successfully corrected by a linear control.

  \subsection*{The networks that encode predictions also encode prediction error}

    Next, we enquired whether the network states \emph{encode} prediction error; i.e., whether the prediction error vector $e_t$ could be decoded through a linear readout of the hidden states. This is a stronger condition than the state update magnitude merely correlating with prediction error.

    To do this, we fitted ridge regressions from the hidden state $h_t$ to the prediction error $e_t = x_t - \hat{x}_{t|t-1}$ using the validation set. We then evaluated the regressions using the 140 compositions of the test set. We contextualised the decoding accuracy using three benchmark models also trained on the validation set: (1)~a constant decoder emitting the expected prediction error vector; (2)~a decoder predicting $e_t$ from the current observation $y_t$; and (3)~a decoder predicting $e_t$ from the current and past observations $y_t$ and $y_{t-1}$. Since the RNN states may encode information about the current and previous observations, we determined that the RNNs encoded prediction error only if they outperformed the third and most stringent of the benchmarks.

    The linear readout of the case-example network outperformed all the benchmarks in all compositions (Figure~\ref{fig:decoding}A; $d = 3.46$, $p < 10^{-20}$ for the comparison against the most stringent benchmark), demonstrating that it encodes prediction error. Robust encoding of prediction error, however, is only apparent for 64-unit networks trained on noise regimes with $0.15 \leq \sigma < 0.6$; i.e., the noise regimes for which the RNNs outperformed the Markovian benchmarks (Figure~\ref{fig:decoding}B). This pattern is replicated for larger RNNs, but prediction error encoding is not apparent in networks with $N < 32$ units.

    \begin{figure}[p]
      \centering
      \includegraphics[width=\textwidth]{./graph/fig06_pe-decoding.pdf}
      \caption{\textbf{RNNs trained at moderate noise levels encode prediction error.}
      \textbf{A)}~Mean squared error (MSE; in log scale) of the prediction error decoded from the RNN hidden states for the case-example RNN ($\sigma = 0.25$, $N = 64$; blue) in comparison with three benchmark models: a constant (dark gray), a linear regressor using as input the current observation (gray), and a linear regressor using as input the current and past observations (light gray). The linear readout outperforms all the benchmarks in all compositions.
      \textbf{B)}~Cohen's $d$ for the contrast between the RNN readout and the best-informed benchmark, which takes as input the current and past observations, across network sizes ($y$-axis) and noise levels ($x$-axis). Values of Cohen's $d$ enclosed in parentheses mark results that were not significant after Bonferroni correction for 84 comparisons.}
      \label{fig:decoding}
    \end{figure}

